% s-167-p167-m003; original PDF 167, 168
\Aphorism{s-167-p167-m003}{256}
Geistige Menschen haben die feinste Haut und den längsten Nachhall für Erlebnisse der äusseren, ungeistigen Sphäre, während gerade eine intellectuelle Wirkung bei ihnen rasch verklingt und leicht auszulöschen ist. Umgekehrt kann auf praktische Menschen, die ohne grosse Erregung ihren Knäul von Geschäften, Peinlichkeiten, selbst Gefahren zu entwirren wissen, nichts tiefer, einschlagender wirken als eine Musik, eine Philosophie: denn hier sind sie fremd, ungeübt, schwerfällig organisirt, wie Jene im Praktischen. — Ärzte z. B. glauben sich bisweilen verpflichtet, Nervenkranken gewisse Arten von Lectüre zu verbieten; sie überschätzen eben, als Männer der Praxis und Handlichkeit, den Eindruck von Lectüre, weil sie selbst wenig lesen. Andererseits neigen wir contemplativen Menschen sehr zur Bewunderung der activen, indem wir ebenfalls den Nervenverbrauch beim Thun und Handeln überschätzen — wir, die wir selbst wenig, viel zu wenig thätig sind. Ich glaube, der moderne Europäer ist in der intellectuellen Verfeinerung schon heute so weit, dass ihm jedes äussere Erleben an die Nerven greift; er hat Fieber beim Aufenthaltswechsel, Fieber bei jeder neuen Bekanntschaft, Reisefieber, Umzugsfieber, Lampen-, Duell-, Kanonenfieber und wie diese Errungenschaften einer verweichlichenden Cultur alle heissen mögen. Also bei uns ist in der That alles Praktische mit erheblichem Nervenverbrauch verbunden, aber wir irren, wenn wir uns die handelnden Menschen in einem dauernden Zustand solcher fieberhaften Spannung und Kraftausgabe vorstellen.
\par

